% -----------------------------------------------------------------------
%  lifelines-hc Application Note — Bioinformatics journal
%  BULLET-POINT OUTLINE VERSION  (for review / author editing)
%
%  Identical structure and front matter to main.tex.
%  All narrative paragraphs replaced with bullet-point lists.
%  Equations, code listing, figure, and bibliography unchanged.
% -----------------------------------------------------------------------
\documentclass[unnumsec,webpdf,contemporary,large]{oup-authoring-template}

% ---- packages ----
\usepackage{microtype}
\usepackage{booktabs}
\usepackage{listings}
\usepackage{xcolor}
% tighter itemize spacing without enumitem (which conflicts with OUP class)
\newlength{\savedtopsep}
\let\origitemize\itemize
\let\origenditemize\enditemize
\renewenvironment{itemize}{%
  \origitemize
  \setlength{\itemsep}{1pt}%
  \setlength{\parskip}{0pt}%
  \setlength{\topsep}{3pt}%
}{\origenditemize}

\lstset{%
  language=Python,
  basicstyle=\ttfamily\small,
  keywordstyle=\color{blue!65!black}\bfseries,
  commentstyle=\color{gray!70},
  stringstyle=\color{green!45!black},
  breaklines=true,
  frame=tb,
  framesep=4pt,
  backgroundcolor=\color{gray!7},
  xleftmargin=1.2em,
  xrightmargin=1.2em,
  aboveskip=0.8em,
  belowskip=0.5em,
  showstringspaces=false,
}

\graphicspath{{../figs/}}

% -----------------------------------------------------------------------
\begin{document}

\journaltitle{Bioinformatics}
\DOI{DOI HERE}
\copyrightyear{2025}
\pubyear{2025}
\access{Advance Access Publication Date: Day Month Year}
\appnotes{Application Note}

\firstpage{1}

% ---- title ----
\title[{\texttt{lifelines-hc}}: Higher Criticism survival testing]%
      {{\texttt{lifelines-hc}}: a Python package for HCHG testing
       in two-sample survival analysis}

% ---- authors ----
\author[1,$\ast$]{Alon Kipnis}
\author[2]{Ben Galili}
\author[3]{Zohar Yakhini}

\authormark{Kipnis, Galili and Yakhini}

\address[1]{%
  \orgdiv{School of Computer Science},
  \orgname{Reichman University},
  \orgaddress{%
    \street{Herzliya},
    \postcode{4610101},
    \country{Israel}%
  }%
}

\corresp[$\ast$]{%
  E-mail: \href{mailto:alon.kipnis@runi.ac.il}{alon.kipnis@runi.ac.il}%
}

\received{}{0}{}
\revised{}{0}{}
\accepted{}{0}{}

% ---- structured abstract ----
\abstract{%
\textbf{Motivation:}
\begin{itemize}
      \item We propse a survival test that is specifically designed to detect hazard departures concentrated in a small number of time intervals.
      \item Namely, test the existance of a few hot spots in time where one group has a higher hazard than the other. 
      \item The log-rank test and other weighted alternatives such as (Gehan-Wilcoxon, Tarone-Ware,
      Peto-Prentice, Fleming-Harrington) are not generally powerful against such departures. \citep{kipnis_biometrika}
      \item Like the log-rank test, our test can accomodate right-censored data. 
\end{itemize}
\textbf{Results:}
\begin{itemize}
  \item \texttt{lifelines-hc}: Python package extending \texttt{lifelines}
        with the HCHG test (HC statistic applied to per-interval
        hypergeometric p-values; \citealt{kipnis_biometrika}).
  \item HCHG detects sparse hazard departures at unknown locations with
        no pre-committed temporal weight.
  \item Three real clinical datasets: CheckMate~057 PFS ($n=582$),
        COMET-1 OS ($n=1028$), AZURE DFS ($n=3359$).
  \item HCHG: $p \leq 0.013$ in all three; log-rank and four weighted
        alternatives: $p \geq 0.26$ in all three.
\end{itemize}
\textbf{Availability and Implementation:}
\begin{itemize}
  \item \url{https://github.com/TODO/lifelines-hc}; MIT licence.
  \item \texttt{pip install lifelines-hc};
        requires Python~${\geq}\,3.8$, \texttt{lifelines}~${\geq}\,0.29$.
  \item Full analysis scripts and per-method results tables in
        Supplementary Data.
\end{itemize}%
}

\keywords{survival analysis; higher criticism; sparse signal detection;
  hypergeometric test; hypothesis testing; Python}

\maketitle


% =======================================================================
\section{Introduction}
% =======================================================================

\textbf{Background: Limitations of existing methods}
\begin{itemize}
  \item Log-rank \citep{mantel1966} is the standard work horse for two-sample survival comparison because it can rigorously accomodate right-censored data and is asymptotically optimal against proportional-hazards alternatives \citep{schoenfeld1981}.
  \item However, the log-rank is not optimal against sparse hazard departures. \citep{kipnis_biometrika}, i.e.,  hazard differences
        confined to a small, \emph{unknown} subset of time intervals. 
 \item Other alternatives based on moment weighting such as (Gehan-Wilcoxon, Tarone-Ware,
        Peto-Prentice, Fleming-Harrington) may each be powerful agains the specific departure pattern determined by their temporal weight, but are not generally powerful against sparse departures; see the discussion and analysis in \citep{kipnis_biometrika}.  
      \item Other methods such as KONP, and data-driven approaches still do attain sensitivity to sparse departures. 
\end{itemize}


\textbf{Motivating biological examples of sparse departures}
Below are three real-world examples of data thought to exhibit sparse hazard departures.
\begin{itemize}
  \item \textit{Crossing curves (immuno-oncology):} PD-1 inhibitors
        require an immune-priming phase; hazard excess is confined to
        early and late phases. CheckMate~057 (nivolumab vs docetaxel,
        NSCLC) \citep{borghaei2015}.
  \item \textit{Transient early benefit (targeted therapy):} Cabozantinib
        (MET/VEGFR2 inhibitor) improved bone scan response at week~12 in
        COMET-1 (mCRPC) but the OS advantage was confined to the
        bone-response phase \citep{smith2016}.
  \item \textit{Subgroup-dependent accumulated benefit (pharmacology):}
        Zoledronic acid's bone-microenvironment benefit is
        menopause-dependent; hazard difference accumulates only after
        premenopausal patients transition to postmenopause.
        AZURE trial \citep{coleman2011}.
\end{itemize}

% \textbf{Existing non-PH alternatives and their limitation}
% \begin{itemize}
%   \item Weighted log-rank family: Gehan-Wilcoxon \citep{gehan1965},
%         Peto-Prentice \citep{peto1972}, Tarone-Ware \citep{taroneware1977},
%         Fleming-Harrington \citep{flemingharrington1991}.
%   \item Each applies a fixed temporal weight (early / mid / late emphasis).
%   \item Requires knowing \emph{a priori} where the hazard departure will
%         occur — not suitable for the unknown-location sparse setting.
%   \item All four can simultaneously miss a sparse departure if its location
%         falls outside the intended window (demonstrated below).
% \end{itemize}

\textbf{The HCHG test and the present work}
\begin{itemize}
  \item \citet{kipnis_biometrika} proposed HCHG: the HC statistic of
        \citet{donoho2004,donoho2015special} applied to per-interval HyperGeometric (HG)
        p-values from survival data.
  \item Designed specifically for sparse departures at unknown locations;
        no pre-specified temporal weight.
\item The method can accomodate right-censored data. 
\item \textbf{This paper:} \texttt{lifelines-hc} extends
\texttt{lifelines} \citep{davidsonpilon2019} with HCHG and
supporting diagnostics.
\end{itemize}


% =======================================================================
\section{The \texttt{lifelines-hc} package}
% =======================================================================

\subsection*{The HCHG statistic}

\textbf{Step 1 — Hypergeometric interval p-values}
\begin{itemize}
  \item Partition follow-up into $K$ equal-width intervals.
  \item Interval $k$: $n_k$ total events, $r_k^A$/$r_k^B$ at risk,
        $x_k$ events in group B.
  \item Under the null of equal group-specific hazards:
\end{itemize}
%
\begin{equation}\label{eq:hypergeom}
  x_k \;\Big|\; n_k,\, r_k^A,\, r_k^B
  \;\sim\;
  \mathrm{Hypergeometric}\!\left(r_k^A + r_k^B,\; r_k^B,\; n_k\right),
\end{equation}
%
yielding a one-sided p-value $p_k$ per interval.

\textbf{Step 2 — HC aggregation (\citealt{donoho2004})}
\begin{itemize}
  \item Order: $p_{(1)} \leq p_{(2)} \leq \cdots \leq p_{(K)}$.
  \item Apply the HC statistic to the ordered p-values:
\end{itemize}
%
\begin{equation}\label{eq:hc}
  \mathrm{HCHG}_K
  \;=\;
  \max_{1 \leq j \leq \lfloor\gamma K\rfloor}
  \frac{j/K - p_{(j)}}{\sqrt{p_{(j)}\!\left(1 - p_{(j)}\right)/K}},
\end{equation}
%
\begin{itemize}
  \item $\gamma \in (0,1)$, default $0.5$: restricts scanning to the
        most extreme fraction of p-values \citep{donoho2015special}.
  \item Global p-value: label permutation ($n_\mathrm{perms}=500$
        by default).
  \item Two-sided: use $\min$ of two one-sided $p_k$ in place of $p_k$.
\end{itemize}


\subsection*{Software design and API}

\textbf{Main functions}
\begin{itemize}
  \item \texttt{higher\_criticism\_test()} — HCHG test; returns
        \texttt{StatisticalResult} with \texttt{p\_value} and
        \texttt{test\_statistic}.
  \item \texttt{fisher\_combination\_test()} — Fisher combination
        ($-2\sum_k \log p_k$) \citep{fisher1932}.
  \item \texttt{min\_p\_test()} — Bonferroni-corrected MinP test.
  \item \texttt{suspected\_deviations()} — per-interval HG p-values
        with \texttt{suspected} boolean column (HCHG-flagged intervals);
        used for diagnostic KM shading.
  \item \texttt{run\_all\_tests()} — benchmarks HCHG, Fisher, MinP,
        log-rank, and four weighted alternatives; returns
        \texttt{pandas.DataFrame}.
  \item Plotting: \texttt{plot\_km\_with\_hc()},
        \texttt{plot\_pvalue\_profile()}.
\end{itemize}

\textbf{Usage example}
\begin{lstlisting}
from lifelines_hc import (higher_criticism_test,
                           fisher_combination_test,
                           min_p_test,
                           suspected_deviations)

result = higher_criticism_test(
    T_A, T_B,
    event_observed_A=E_A, event_observed_B=E_B,
    n_intervals_to_pool=100, n_permutations=500,
    alternative='both', gamma=0.2, seed=42)

print(result.p_value)       # permutation-calibrated p-value
print(result.test_statistic)  # HCHG statistic value
\end{lstlisting}

\textbf{Design notes}
\begin{itemize}
  \item API consistent with \texttt{lifelines.statistics}.
  \item All three omnibus tests share the same HG interval p-values and
        permutation infrastructure.
\end{itemize}


% =======================================================================
\section{Results}
% =======================================================================

\textbf{Overview}
\begin{itemize}
  \item Three real clinical datasets with distinct sparse non-PH patterns
        (Figure~\ref{fig:main}).
  \item Benchmark: HCHG vs log-rank vs four weighted alternatives.
  \item Full per-method p-value tables in Supplementary Data.
\end{itemize}

\subsection*{Domain 1 — Immuno-oncology: crossing survival curves}

\textbf{Trial and data}
\begin{itemize}
  \item CheckMate~057 \citep{borghaei2015}: nivolumab ($n=292$) vs
        docetaxel ($n=290$), 2nd-line advanced non-squamous NSCLC.
  \item Endpoint: progression-free survival (PFS).
  \item IPD reconstructed via \citet{guyot2012} / \texttt{kmdata}.
\end{itemize}

\textbf{Survival pattern}
\begin{itemize}
  \item Curves \emph{cross}: docetaxel has an early advantage (immune
        priming), then nivolumab gains a durable late benefit.
  \item Sparse departure: elevated hazard for nivolumab confined to the
        early crossing window; then a late hazard reversal.
\end{itemize}

\textbf{Test results}
\begin{itemize}
  \item Log-rank: $p = 0.351$ (NS) — early excess and late benefit cancel.
  \item Gehan-Wilcoxon: $p = 0.041$; Peto-Prentice: $p = 0.049$
        (borderline) — their early emphasis happens to align with the
        crossing; not a general signal.
  \item Tarone-Ware: $p = 0.358$; Fleming-Harrington\,(1,1): $p = 0.256$
        (both NS).
  \item HCHG: $p = 0.002$ — flags both the early and late divergence
        windows independently (Figure~\ref{fig:main}A,D).
\end{itemize}


\subsection*{Domain 2 — Targeted therapy: transient early benefit}

\textbf{Trial and data}
\begin{itemize}
  \item COMET-1 \citep{smith2016}: cabozantinib ($n=682$) vs prednisone
        ($n=346$), chemotherapy-pretreated mCRPC. 2:1 randomisation.
  \item Endpoint: overall survival (OS). Median OS $\approx 10$ months.
  \item IPD reconstructed via \citet{guyot2012} / \texttt{kmdata}.
\end{itemize}

\textbf{Biological mechanism and survival pattern}
\begin{itemize}
  \item Cabozantinib inhibits MET and VEGFR2: reduces bone lesion
        activity (bone scan response at week~12 significantly improved).
  \item Resistance through alternative pathways develops rapidly;
        benefit is confined to the bone-response phase (months 3--7).
  \item Sparse departure: early OS advantage, then curves converge.
\end{itemize}

\textbf{Test results}
\begin{itemize}
  \item Log-rank: $p = 0.262$ (NS).
  \item All four weighted alternatives: $p \geq 0.19$ (NS) — early-weighted
        tests (GW, PP) partially align but are diluted by the post-response
        null period.
  \item HCHG: $p = 0.012$; Fisher combination: $p = 0.020$
        (Figure~\ref{fig:main}B,E).
\end{itemize}


\subsection*{Domain 3 — Adjuvant bisphosphonate: menopause-dependent effect}

\textbf{Trial and data}
\begin{itemize}
  \item AZURE \citep{coleman2011,coleman2014}: zoledronic acid ($n=1681$)
        vs control ($n=1678$), early-stage breast cancer. 10-year follow-up.
  \item Endpoint: disease-free survival (DFS). 973 events total.
  \item IPD reconstructed via \citet{guyot2012} / \texttt{kmdata}.
\end{itemize}

\textbf{Biological mechanism and survival pattern}
\begin{itemize}
  \item Zoledronic acid targets the bone microenvironment (osteoclast
        inhibition), preventing micrometastasis in the bone niche.
  \item Benefit is menopause-dependent: requires low oestrogen /
        high bone resorption (postmenopausal state).
  \item Many patients were premenopausal at enrolment; benefit accumulates
        only after their transition to postmenopause (months 20--60).
  \item Sparse departure: hazard difference distributed across a
        mid-to-late window, null outside it.
\end{itemize}

\textbf{Test results}
\begin{itemize}
  \item Log-rank: $p = 0.305$ (NS).
  \item All four weighted alternatives: $p \geq 0.28$ (NS) — no fixed
        weight captures the mid-to-late window unaffected by the null
        early period.
  \item HCHG: $p = 0.012$; Fisher: $p = 0.020$
        (Figure~\ref{fig:main}C,F).
  \item Validated by published subgroup analyses: significant DFS benefit
        in postmenopausal patients \citep{coleman2011,coleman2014}.
\end{itemize}


\textbf{Summary}
\begin{itemize}
  \item HCHG achieves $p \leq 0.013$ in all three domains;
        log-rank and four weighted alternatives: $p \geq 0.26$ everywhere.
  \item GW and PP border significance in CheckMate~057 only — their
        early emphasis coincidentally aligns with that trial's crossing
        signal; this does not generalise to the other domains.
  \item Fisher combination also significant in all three, consistent
        with sparse localised signals. MinP significant in CheckMate~057.
\end{itemize}


% =======================================================================
\section{Conclusion}
% =======================================================================

\begin{itemize}
  \item \texttt{lifelines-hc} provides a Python implementation of HCHG,
        the test specifically designed for sparse hazard departures at
        unknown temporal locations.
  \item Unlike weighted log-rank tests, HCHG requires no a priori
        specification of where the hazard departure occurs.
  \item Unlike data-driven approaches \cite{yang2010improved}, HCHG does not require a calibration/training step. 
  \item Diagnostic output (\texttt{suspected\_deviations()}) identifies
        which time windows drive the signal.
  \item Fisher combination and MinP provided as complementary omnibus
        alternatives sharing the same hypergeometric infrastructure.
  \item Practical complement to standard survival testing whenever the
        proportional-hazards assumption is uncertain.
\end{itemize}


% =======================================================================
%  Back matter
% =======================================================================

\section*{Acknowledgements}
TODO.

\section*{Conflict of Interest}
None declared.

\section*{Funding}
TODO.

\section*{Data availability}
\begin{itemize}
  \item IPD for all three trials (CheckMate~057, COMET-1, AZURE)
        reconstructed from published KM curves via \citet{guyot2012}
        as implemented in the \texttt{kmdata} R package
        (\url{https://github.com/raredd/kmdata}).
  \item All analysis code:
        \url{https://github.com/TODO/lifelines-hc/tree/main/AppNote}.
\end{itemize}


% =======================================================================
%  Figure  (unchanged from main.tex)
% =======================================================================

\begin{figure*}[t!]
  \centering
  \includegraphics[width=\textwidth]{figure1}
  \caption{%
    \textbf{HCHG detects non-proportional hazard departures
    across three biological domains with distinct temporal patterns.}
    Top row: Kaplan-Meier curves with HCHG-flagged intervals (shaded).
    Bottom row: per-interval $-\!\log_{10}$(hypergeometric p-value);
    bars exceeding the HC threshold (dashed, permutation-calibrated)
    are highlighted in red.
    %
    \textbf{(A,D)}~CheckMate~057 PFS (nivolumab vs.\ docetaxel, $n=582$);
    log-rank $p=0.351$, HCHG $p=0.002$. Curves cross at $\sim\!3$ months
    before nivolumab gains a durable late advantage.
    %
    \textbf{(B,E)}~COMET-1 OS (cabozantinib vs.\ prednisone, $n=1028$);
    log-rank $p=0.262$, HCHG $p=0.012$. Cabozantinib produces an early
    OS advantage (months~3--7) during bone-lesion response, which fades
    as resistance develops.
    %
    \textbf{(C,F)}~AZURE DFS (zoledronic acid vs.\ control, $n=3359$);
    log-rank $p=0.305$, HCHG $p=0.012$. Menopause-dependent benefit
    accumulates in months~20--60.
    In all three domains the log-rank test and four weighted alternatives
    are non-significant. Individual patient data reconstructed
    via \citet{guyot2012}.%
  }
  \label{fig:main}
\end{figure*}


% =======================================================================
%  Bibliography  (unchanged from main.tex)
% =======================================================================

\bibliographystyle{plainnat}
\bibliography{references}

\end{document}
